El panorama descrito evidencia una problemática estructural en la enseñanza de la trigonometría en educación media colombiana. Por una parte, persisten prácticas centradas en memorización y resolución rutinaria que limitan la comprensión funcional de seno, coseno y tangente \parencite{fiallolealAcercaInvestigacionEducacion2015,delgadoEstrategiaDidacticaPara2023}. Por otra, los estudiantes enfrentan obstáculos conceptuales recurrentes para interpretar la variación periódica, coordinar diferentes representaciones y argumentar matemáticamente sus decisiones \parencite{zubietaTipificacionErroresDificultades2018,ruedaAproximacionEnsenanzaRazones2012}.

Esta situación genera una brecha entre las metas curriculares de educación media y lo que efectivamente ocurre en el aula. Aunque los lineamientos nacionales enfatizan pensamiento variacional, modelación y comprensión de funciones, en la práctica muchos procesos de enseñanza no logran sostener ese tránsito desde la razón trigonométrica hacia el pensamiento trigonométrico \parencite{ministeriodeeducacionnacionalLineamientosCurriculares1998,ministeriodeeducacionnacionalEstandaresBasicos2006,ministeriodeeducacionnacionalDerechosBasicos2016}.

En este contexto, la integración con inteligencia artificial no se asume como un fin en sí mismo ni como sustituto de la actividad matemática escolar. El reto está en diseñar e implementar tareas matemáticas integradas con inteligencia artificial (IA) para apoyar procesos de exploración, argumentación y retroalimentación, manteniendo como foco el desarrollo del pensamiento trigonométrico en estudiantes de educación media \parencite{castillodelpezoUseArtificialIntelligence2024,parrasanchezPersonalizacionRecursosPara2023}.
